\section{What to listen for at Mass and in the Hours}

\subsection{The Sunday collect}

When the priest invites us to pray, there is time to bring an intention to God. The collect gathers the assembly into a single petition. The four Sunday collects of Advent return each year, whichever cycle of Gospels the parish is reading. A child who learns their movement can hear a new Gospel against a familiar prayer.

On the First Sunday, listen for the people running and for the place at which they hope to arrive. They ask for the will to go out to Christ with good works, and for a place at his right hand. In Matthew 25:31--40 the King at whose right hand the just stand remembers a meal, a garment, a visit made to a prisoner. It is a startling choice of image for the beginning of a season when much of our running about is already spoken for. Who could use the hour between the last errand and supper? Notice who supplies even the will to set out. The priest asks the Father for it before any of us has done a thing.

On the Second Sunday the prayer speaks of hindrances and asks for heavenly wisdom. It is possible to have done good work all morning and to resent the first person who interrupts it. Wisdom would let us see what this particular duty, this particular interruption, asks of us. John will have much to say about preparing a straight road. In the collect we ask to learn the way in company with the Son. On the Third Sunday the petition itself is for joy. Perhaps gladness is the one thing you feel unable to give. At Gaudete Mass you can ask for it with the rest of the parish. And if the person beside you is grieving, let your presence there spare him the burden of having to appear cheerful. The rose vestments which may be worn take their sense from the day's invitation to rejoice in the Lord (Philippians 4:4--5).

On the Fourth Sunday the priest recalls the angel's message and asks that we be brought through the Passion and Cross to the Resurrection. Those who pray the Angelus will recognize its closing prayer. A mother has scarcely heard that she will bear a son, and the Church is already praying about Easter. The name Jesus, which Joseph is told to give the child, means something about what he will do for his people (Matthew 1:21). We have come to Bethlehem to welcome our Redeemer. Paul VI drew attention to this very prayer when he commended the Angelus to the faithful (\emph{Marialis cultus}, no.~41).

\subsection{From the ambo to the altar}

The reader sits down and the psalmist sings. Before hurrying on to the Gospel, try praying the response in the name of someone else. A plea for deliverance can belong to a Christian far from home, or to a neighbour who has grown weary of asking. The psalm belongs to the day's proclamation of Scripture. In making its words our own, we learn to answer the God who has just spoken (General Instruction of the Roman Missal, no.~61). Advent Sundays retain the Alleluia. The Gloria ordinarily waits until the feast, and the Creed is professed on Sunday. The white vestments and the Gloria at a solemnity such as the Immaculate Conception belong to that celebration.

Listen again when the priest begins the great prayer of thanksgiving. Before December 17 the first Advent preface remembers the Lord's lowliness in taking flesh and the glory in which he will return. From December 17 the second brings the prophets, his Mother and John the Baptist into our thanksgiving. The prophets longed for him, Mary bore him and John pointed him out. We come after all three. Their hope is already answered, and our generation still has the privilege of keeping watch. Then our voices join the angels in the Sanctus. At Communion the one whose coming we have praised gives us his Body to eat. To go from that table and neglect those who go hungry would be a grievous way of forgetting whose Body we have received (1 Corinthians 11:17--29).

\subsection{A morning with Zechariah, an evening with Mary}

Morning Prayer and Evening Prayer are the two principal Hours of the Church's daily prayer. A parish can celebrate them together, a family can pray them at home, and someone praying alone can join in the Church's praise. In the U.S. four-volume \emph{Liturgy of the Hours}, Advent is in volume I. Begin with the Proper of Seasons for the day and use the week of the Psalter indicated there. A parish that offers Evening Prayer is a good place to learn when to listen, when to sing and when to stand. The General Instruction on the Liturgy of the Hours, nos.~37--54, sets out these two Hours.

At Morning Prayer the Gospel canticle is the \emph{Benedictus} (Luke 1:68--79). Listen for the moment the father turns to address his child. Zechariah begins with Abraham, David and the oath of God. Only when the whole house of Israel has come into view does he speak to the little boy in front of him. John will go before the Lord, teaching his people where salvation is to be found. On the last morning of Advent that very canticle is the weekday Gospel. A Christian who has said it each morning may hear, for the first time, a father speaking to a newborn.

At Evening Prayer the Gospel canticle is Mary's \emph{Magnificat} (Luke 1:46--55). She speaks of what God has done for her, then of his mercy to generations before and after her. When she names the hungry at the end of a working day, people in your own town are sitting down to eat. Some will be alone. Others will count the portions quite carefully. Remember one of them while praying the canticle. If you have food or company to share, arrange a time to bring it. The Church places intercessions and the Our Father after the canticle. Even a person praying alone can hear the plural in the prayer our Lord taught us.

The Office of Readings gives more time to listen. A biblical lesson is followed, in the ordinary course, by a second reading from a Father or another ecclesiastical writer and by a responsory after each lesson (General Instruction, nos.~55--56, 64--65). Look up the First Sunday of Advent in volume I and read St Cyril of Jerusalem on the two advents. Return on Wednesday of that week for St Bernard's sermon on the coming which takes place within the believer. Someone who wonders how Christ could come today will find in those two readings both the Church's promise and a way of receiving it.

\subsection{Seven names to say before Christmas}

From December 17 through 23 the Church sings an antiphon to Christ with Mary's song at Evening Prayer. The first word of each is an address, \emph{O}. At home you can use the traditional Latin name, read the short passage suggested here, and pray Luke 1:46--55. These passages are companions for prayer. The readings of the weekday Mass are given in their own chapter.

\paragraph{December 17. Wisdom.} \emph{O Sapientia.} Read Wisdom 8:1. Wisdom can order all things gently. Think of a decision that is yours to make for a child or an elderly parent. Ask Christ for the wisdom to do what is good without humiliating anyone who depends upon you.

\paragraph{December 18. Adonai.} \emph{O Adonai.} In Exodus 3:7--10 the Lord tells Moses what he has seen and heard in Egypt. Then he sends Moses to Pharaoh. The title \emph{Adonai} addresses the Lord who hears the cry of his people. Spend a moment listening to a cry that ordinarily reaches you only as a piece of news.

\paragraph{December 19. Root of Jesse.} \emph{O Radix Jesse.} Read Isaiah 11:1, 10. Jesse was David's father. A felled tree looks like the end of a royal house, yet a shoot grows from its roots. The nations will come to the one who springs from that family. Name before Christ a people waiting for peace.

\paragraph{December 20. Key of David.} \emph{O Clavis David.} In Isaiah 22:22 a key is entrusted to the steward of a household. Revelation 3:7 applies the image to Christ. Remember a prisoner, or someone whose past has closed doors even after a sentence has been served. Ask Christ to open the door, and consider who will be waiting to welcome that person when it opens.

\paragraph{December 21. Dayspring.} \emph{O Oriens.} Read Luke 1:78--79. \emph{Oriens} names the rising light. In Zechariah's song the dawn reaches people sitting in darkness and guides their feet. This evening give thanks for the person who stayed beside you on a long night. If you yourself are keeping watch, you can say the last two verses of the canticle before morning comes.

\paragraph{December 22. King of the nations.} \emph{O Rex Gentium.} Read Ephesians 2:13--19. A cornerstone joins walls that approach from different directions. Paul speaks of Jews and Gentiles reconciled in the one Body through the Cross. Bring to Christ a division in which you have a share. Ask what peace would cost you and who has already borne its greatest cost.

\paragraph{December 23. Emmanuel.} \emph{O Emmanuel.} Read Matthew 1:20--24. The angel tells Joseph what to name the child, and the evangelist remembers Isaiah's promise of God with us. Joseph takes Mary into his home. As Christmas approaches, consider who in your household has borne the heavier part of its preparations. Take up one task in time for the other person to rest.

\subsection{Making a good confession}

The crowds who came to John asked what they should do. His answers had to do with clothing, food, honest wages, the use of power (Luke 3:10--14). If it has been a long time since your last confession, read those few verses. Bring to mind the people who have had to live with your decisions. Ask the Holy Spirit for sorrow, and tell the priest plainly what you have done. If you have forgotten how to begin, say so. A priest can help you find the words. Receive the absolution Christ gives through his minister, do the penance assigned, and make good what can be repaired. An apology or the return of money can be part of that work. The \emph{Catechism}, nos.~1451, 1455--1460, explains why forgiven sins may still leave a neighbour with a claim upon us.

There is room for a freely chosen sacrifice in Advent. In the Latin Church the universal penitential days and seasons are Fridays and Lent. In the United States the bishops especially commend abstinence from meat on Fridays outside Lent and allow other forms of penance. A solemnity changes the ordinary Friday observance. Advent brings no general obligation to fast each weekday. Someone feeding children can set aside a portion for another household. Someone who is ill can accept the nourishment needed for recovery and remember the hungry in prayer. For the law and its U.S. application, see the appendix on the books and calendar.
